% =====================================================================
% SZKIC v1 (01.11.2026) -- wersja polska (robocza, do dyskusji z promotorem)
% Odpowiada main.tex (EN). Protokół identyczny z poprzednimi pracami;
% nowa jest tylko klasa sąsiedztw. Tabela 2 = PILOT, Tabela 3 = szkielet.
% =====================================================================
\documentclass[10pt,twocolumn]{article}
\usepackage[margin=0.75in]{geometry}
\usepackage[T1]{fontenc}
\usepackage[utf8]{inputenc}
\usepackage[polish]{babel}
\usepackage{amsmath,amssymb,amsthm}
\usepackage{booktabs}
\usepackage{multirow}
\usepackage{xcolor}
\usepackage[hidelinks]{hyperref}

\newtheorem{proposition}{Twierdzenie}
\newtheorem{definition}{Definicja}
\newcommand{\todo}[1]{\textcolor{red}{[DO ZROBIENIA: #1]}}
\newcommand{\N}{\mathcal{N}}

\title{Od Fibonacciego do Dynasearch: rodzina sąsiedztw $k$-bonacci\\
złożonych z zamian dla permutacyjnego problemu przepływowego}
\author{Zgłoszenie anonimowe}
\date{}

\begin{document}
\maketitle

\begin{abstract}
Sąsiedztwa ruchów złożonych pozwalają przeszukać wykładniczo wiele
uszeregowań w czasie wielomianowym metodą programowania dynamicznego.
Dwa najbardziej znane sąsiedztwa tego typu dla permutacyjnego problemu
przepływowego leżą jednak na przeciwnych biegunach. Sąsiedztwo
Fibonacciego łączy tylko zamiany sąsiednie i jest tanie, natomiast
Dynasearch łączy zamiany na dowolną odległość i kosztuje $O(mn^3)$ na
ocenę. Proponujemy rodzinę \emph{$k$-bonacci}, która wypełnia lukę między
nimi: ruch to zbiór zamian końców na rozłącznych przedziałach, przy
odległości zamiany co najwyżej $k-1$. Liczność rodziny wyznaczają
dokładnie liczby $k$-bonacciego, więc rośnie ona jak $\lambda_k^n$, gdzie
$\lambda_2=1{,}618$ (Fibonacci), $\lambda_3=1{,}839$ (tribonacci)
i $\lambda_k\to 2$ (Dynasearch). Już dla $k\approx\log_2 n+2$ rodzina
zawiera ok. 80\% wszystkich ruchów Dynasearch. Pozycyjne programowanie
dynamiczne wybiera najlepszy ruch w czasie $O(nk)$. Rodzinę badamy na
pełnym benchmarku Taillarda w ramach ILS i SA, przy budżetach czasu
z naszych wcześniejszych prac, dzięki czemu nowe sąsiedztwa można
bezpośrednio porównać z opublikowanymi wynikami dla sąsiedztw Adjacent,
Fibonacciego, Dynasearch i Motzkina.
\todo{główny wynik: najlepsze $k$ dla rozmiaru i budżetu}
\end{abstract}

\section{Wstęp}
Permutacyjny problem przepływowy z kryterium $C_{\max}$,
$F|prmu|C_{\max}$, jest NP-trudny dla $m\geq 3$ maszyn~\cite{garey1976}.
O jakości przeszukiwania lokalnego decydują przede wszystkim dwie rzeczy:
struktura ruchu oraz własności blokowe ścieżki krytycznej, które pozwalają
ten ruch przycinać~\cite{nowicki1996,grabowski2004}.

Dynasearch~\cite{congram2002} pokazał, że wykładniczo duże sąsiedztwo
\emph{niezależnych} zamian można przeszukać dokładnie metodą programowania
dynamicznego. Dla PFSP znane są też sąsiedztwo Fibonacciego, złożone
z rozłącznych zamian sąsiednich~\cite{fibonacci_own,bozejko2026sr},
oraz sąsiedztwo Motzkina~\cite{motzkin_own}. Te sąsiedztwa różnią się
o rzędy wielkości rozmiarem i kosztem, a mimo to porównywano je dotąd
wyłącznie jako odosobnione punkty. Praca wypełnia przestrzeń między
Fibonaccim a Dynasearch jednoparametrową rodziną i pyta: \emph{jak daleko
powinna sięgać zamiana?} Metaheurystyki, ich parametry i protokół
eksperymentu celowo pozostawiamy takie jak w poprzednich
pracach~\cite{fibonacci_own,motzkin_own,windowed_own}. Jedynym czynnikiem
eksperymentalnym jest sąsiedztwo, więc każdą nową liczbę można zestawić
z wynikami już opublikowanymi.

\paragraph{Wkład.}
\begin{enumerate}
  \item Sąsiedztwo $k$-bonacci $\N_k$ z łańcuchem zawierań
        $\N_{\mathrm{adj}}\subset\N_2=\N_{\mathrm{fib}}\subset\N_3\subset
        \dots\subset\N_n=\N_{\mathrm{dyn}}$. Podajemy dokładną liczność
        i pokazujemy, że $k=\log_2 n+O(1)$ pokrywa stały ułamek Dynasearch.
  \item Pozycyjne DP o złożoności $O(nk)$ po wyznaczeniu delt w czasie
        $O(mnk^2)$, zgodne z filtrem blokowym. Dla $k=n$ jest to
        nierekurencyjny Dynasearch.
  \item Studium parametru $k$ na pełnym benchmarku Taillarda
        ($n=20$--$500$) w ramach ILS i SA, przy sześciu budżetach czasu.
        Obejmuje jakość, liczbę iteracji, strukturę ruchów i błąd
        addytywności.
  \item \todo{opcjonalnie: QUBO $k$-bonacciego na wyżarzaczu kwantowym
        i komputerze bramkowym, protokół jak w~\cite{windowed_own}}
\end{enumerate}

\section{Prace pokrewne}
Dynasearch wprowadzono dla problemu jednomaszynowego~\cite{congram2002}.
Ergun i Orlin~\cite{ergun2006} przyspieszyli przeszukiwanie do $O(n^2)$,
a Grosso i in.~\cite{grosso2004} dodali wstawienia. W tym problemie zmiany
kosztu niezależnych zamian się sumują, a w PFSP nie
(rozdz.~\ref{sec:addytywnosc}). Ograniczanie zasięgu ruchu to znane
narzędzie w problemach trasowania: Or-opt i sąsiedztwa
granularne~\cite{toth2003}. \todo{sprawdzić, czy Sci Rep 2026 i IFAC 2025
nie zawierają już tej hierarchii}

\section{Problem i notacja}
Mamy $n$ zadań wykonywanych na maszynach $1,\dots,m$ w tej samej
kolejności $\pi$. Czasy zakończenia spełniają
$C_{r,t}=\max(C_{r-1,t},C_{r,t-1})+p_{r,\pi(t)}$. Macierze Head/Tail
pozwalają ocenić modyfikację odcinka $[i,j]$ w czasie $O(m(j-i+1))$.

\section{Sąsiedztwo $k$-bonacci}
\begin{definition}
Niech $k\geq 2$. \emph{Ruch elementarny} $s(i,j)$, gdzie $1\leq j-i\leq k-1$,
zamienia $\pi(i)$ z $\pi(j)$, a zadania $\pi(i+1),\dots,\pi(j-1)$ zostają
na swoich miejscach. \emph{Ruch złożony} to zbiór ruchów elementarnych
o parami rozłącznych przedziałach $[i,j]$.
\end{definition}
Przypadki szczególne: dla $k=2$ dopuszczalna jest tylko odległość $1$
(Fibonacci), dla $k=3$ odległości $1$ i $2$ (tribonacci), a dla $k=n$
dowolna odległość (Dynasearch).

\begin{proposition}
$|\N_k|=a_n$, gdzie $a_0=1$ oraz $a_p=\sum_{L=1}^{\min(k,p)}a_{p-L}$.
Stąd $|\N_2|=F_{n+1}$, $|\N_n|=2^{n-1}$ i $|\N_k|=\Theta(\lambda_k^n)$,
przy czym $2-\lambda_k=\lambda_k^{-k}$, czyli $\lambda_k=2-\Theta(2^{-k})$.
\end{proposition}
\begin{proof}[Szkic]
Ruch złożony to pokrycie pozycji kafelkami długości $1$ (pozycja bez
zmian) lub $L\in\{2,\dots,k\}$ (zamiana końców kafelka). Kafelek
zaczynający się w $i$ to najmniejszy przedział od $i$, który ruch
przeprowadza na siebie, więc różne pokrycia dają różne permutacje.
\end{proof}

\begin{table}[t]
\centering
\caption{Liczność $\N_k$ dla $n=10$, stała wzrostu i udział ruchów
Dynasearch zawartych w $\N_k$.}
\small
\begin{tabular}{@{}lrrrrrr@{}}
\toprule
$k$ & 2 & 3 & 4 & 6 & 8 & $n$ \\
\midrule
$|\N_k|$, $n=10$ & 89 & 274 & 401 & 492 & 509 & 512 \\
$\lambda_k$ & 1,618 & 1,839 & 1,928 & 1,984 & 1,996 & 2 \\
udział, $n=100$ & $10^{-9}$ & $3{\cdot}10^{-4}$ & 0,03 & 0,46 & 0,83 & 1 \\
udział, $n=500$ & $10^{-46}$ & $10^{-18}$ & $10^{-8}$ & 0,02 & 0,38 & 1 \\
\bottomrule
\end{tabular}
\end{table}

Udział $|\N_k|/|\N_n|$ to prawdopodobieństwo, że najdłuższa część losowej
kompozycji $n$ nie przekracza $k$. Zgodnie z prawem najdłuższej serii jest
ono skupione wokół $k=\log_2 n$. Dla $k=\lfloor\log_2 n\rfloor+2$ udział
wynosi 0,79--0,89 przy $n=20$--$500$.

\paragraph{Ocena i DP.} Wszystkie delty kosztują $O(mnk^2)$, a w przypadku
Dynasearch $O(mn^3)$. Kandydatów filtrujemy tym samym testem granic bloków,
którego użyto we wcześniejszych pracach~\cite{windowed_own}.
\todo{decyzja o wariancie dokładnym (plan, \S14)} Pozycyjne DP ma postać
$B[p]=\min\bigl(B[p-1],\ \min_{d} B[p-d-1]+\delta_{p-d,p}\bigr)$ i złożoność
$O(nk)$. Gdy żaden ruch nie poprawia przewidywanego $C_{\max}$, wykonujemy
najlepszą pojedynczą zamianę, tak jak w Dynasearch i Motzkinie.

\section{Addytywność}\label{sec:addytywnosc}
DP minimalizuje $\sum\delta_{ij}$, jednak w PFSP zmiana $C_{\max}$ pod
wpływem ruchu złożonego nie jest sumą zmian cząstkowych. W ujęciu ścieżek
kratowych $C'_{\max}=\max_P[\ell(P)+\sum_s\Delta_s(P)]$, więc zamiany
oddziałują przez wybór ścieżki krytycznej. Każde sąsiedztwo złożone oparte
na DP optymalizuje zatem zastępczą funkcję celu.

\section{Protokół eksperymentu}
Protokół jest taki sam jak we wcześniejszych pracach~\cite{windowed_own}.
Nowe jest tylko sąsiedztwo.
\begin{itemize}
\item \textbf{Instancje:} 120 instancji Taillarda (12 klas,
      $n\in\{20,\dots,500\}$, $m\in\{5,10,20\}$), zweryfikowanych
      generatorem Taillarda.
\item \textbf{Metaheurystyki:} ILS z listą tabu (kadencja 10) i pulą
      elitarną 10 rozwiązań z kickiem double-bridge. SA z $T_0=1000$,
      $\alpha=0{,}99$, podłogą temperatury 100 i reheatem $\times1{,}5$
      z kickiem po 100\,ms stagnacji. Start z permutacji identycznościowej.
\item \textbf{Budżety:} $\{0{,}1;\,0{,}5;\,1;\,2;\,5;\,10\}$\,s,
      5 ziaren (0--4).
\item \textbf{Sąsiedztwa:} $k\in\{2,3,4,8,n\}$. Wiersze Adjacent,
      Fibonacci, Dynasearch i Motzkin pochodzą z wcześniejszej kampanii
      na tej samej maszynie.
\item \textbf{Miara:} RPD względem dolnego ograniczenia Taillarda; test
      Wilcoxona, $\alpha=0{,}05$.
\item \textbf{Sprzęt:} Apple M4 Pro, przebiegi sekwencyjne, jak wcześniej.
\end{itemize}
\paragraph{Jeden przebieg na rodzinę budżetów.} ILS i SA używają limitu
czasu wyłącznie jako warunku stopu. Przebieg z budżetem $t$ jest więc
prefiksem przebiegu 10\,s z tym samym ziarnem. Każdą konfigurację
uruchamiamy raz na 10\,s, a wynik dla mniejszych budżetów odczytujemy
z logu trajektorii. Ruch rozpoczęty przed $t$ jest liczony, jak w niezależnym
przebiegu. \todo{walidacja na dwóch klasach}

\section{Wyniki}
\begin{table}[t]
\centering
\caption{PILOT (1\,s, RPD do UB, 4 instancje $\times$ 3 ziarna):
średnie RPD [\%] / liczba iteracji ILS.}
\small
\begin{tabular}{@{}lccc@{}}
\toprule
Sąsiedztwo & $20\times5$ & $20\times10$ & $50\times10$ \\
\midrule
Adjacent         & 1,59 / 4578 & 3,26 / 2246 & 11,73 / 403 \\
$k=2$ (Fibonacci) & 1,31 / 12400 & 2,46 / 7129 & 8,37 / 3329 \\
$k=3$ (tribonacci) & 0,92 / 10177 & 1,84 / 5007 & 5,94 / 2481 \\
$k=4$            & 0,89 / 7897 & 1,43 / 3618 & 5,51 / 1966 \\
$k=6$            & 0,74 / 5397 & 1,09 / 2375 & 4,85 / 1123 \\
$k=8$            & 0,72 / 3893 & 1,35 / 1580 & 4,66 / 809 \\
$k=n$            & 0,55 / 1539 & 0,99 / 610 & 3,62 / 63 \\
\bottomrule
\end{tabular}
\end{table}

\begin{table*}[t]
\centering
\caption{SZKIELET: średnie RPD [\%] do LB, ILS, dla każdego $n$ (średnia
po $m$, 10 instancji, 5 ziaren). Analogiczna tabela dla SA.}
\small
\begin{tabular}{@{}llcccccc@{}}
\toprule
$n$ & Sąsiedztwo & 0,1\,s & 0,5\,s & 1\,s & 2\,s & 5\,s & 10\,s \\
\midrule
20--500 & Adjacent (ref.) & -- & -- & -- & -- & -- & -- \\
& Fibonacci (ref.) & -- & -- & -- & -- & -- & -- \\
& Dynasearch (ref.) & -- & -- & -- & -- & -- & -- \\
& Motzkin (ref.) & -- & -- & -- & -- & -- & -- \\
\cmidrule(l){2-8}
& $k=2$ & -- & -- & -- & -- & -- & -- \\
& $k=3$ (tribonacci) & -- & -- & -- & -- & -- & -- \\
& $k=4$ & -- & -- & -- & -- & -- & -- \\
& $k=8$ & -- & -- & -- & -- & -- & -- \\
& $k=n$ & -- & -- & -- & -- & -- & -- \\
\bottomrule
\end{tabular}
\end{table*}

\paragraph{Pilot.} Jakość poprawia się monotonicznie od $k=2$ do $k=6$,
a przy 1\,s pełne $k=n$ nadal wygrywa. Błąd addytywności rośnie z $k$:
predykcja jest dokładna w 42\% ocen dla $k=2$ i w 10--12\% dla $k\geq 6$.
\todo{zastąpić wynikami pełnej kampanii}

\section{Wnioski}
\todo{po pełnej kampanii}

\begin{thebibliography}{9}\small
\bibitem{bozejko2026sr} W.~Bożejko, M.~Uchroński, M.~Wodecki, Sci.\ Rep.\ 2026.
\bibitem{congram2002} R.K.~Congram, C.N.~Potts, S.L.~van de Velde, INFORMS J.\ Comput.\ 14(1), 2002.
\bibitem{ergun2006} Ö.~Ergun, J.B.~Orlin, Oper.\ Res.\ Lett.\ 34(1), 2006.
\bibitem{fibonacci_own} \todo{praca własna -- Fibonacci}
\bibitem{garey1976} M.R.~Garey, D.S.~Johnson, R.~Sethi, Math.\ Oper.\ Res.\ 1(2), 1976.
\bibitem{grabowski2004} J.~Grabowski, M.~Wodecki, Comput.\ Oper.\ Res.\ 31(11), 2004.
\bibitem{grosso2004} A.~Grosso, F.~Della Croce, R.~Tadei, Oper.\ Res.\ Lett.\ 32(1), 2004.
\bibitem{motzkin_own} \todo{praca własna -- Motzkin, SOCO 2026}
\bibitem{nowicki1996} E.~Nowicki, C.~Smutnicki, EJOR 91(1), 1996.
\bibitem{toth2003} P.~Toth, D.~Vigo, INFORMS J.\ Comput.\ 15(4), 2003.
\bibitem{windowed_own} \todo{praca własna -- windowed QUBO}
\end{thebibliography}
\end{document}
